% ============================================================
%  Chapter 5 – Results and Discussion
%  Thesis: Modern Control Methods for Agricultural Irrigation
%  Author: Tara Torbati — ITMO University, R4237c
%  Supervisor: Peregudin A. A.
%
%  PLACEHOLDER convention
%  \TODO{description} marks every value / figure that requires
%  a pending experiment to be completed before final submission.
%  Grep for \TODO before submitting.
% ============================================================

\newcommand{\TODO}[1]{\textbf{\textcolor{red}{[TODO: #1]}}}

\chapter{Results and Discussion}
\label{ch:results}

% ============================================================
\section{Introduction}
\label{sec:results_intro}
% ============================================================

This chapter presents the complete empirical evaluation of the four control strategies
developed and described in Chapters~3 and~4: the no-irrigation baseline, the fixed-schedule
heuristic, the Model Predictive Controller (MPC), and the Soft Actor-Critic Reinforcement
Learning agent (SAC).
All controllers are evaluated on the same validated Agent-Based Model (ABM) of the
6-hectare Hashemi rice field in Gilan Province, Iran, across three climatically distinct
test years — dry (2022, 39.7~mm seasonal rainfall), moderate (2018, 108.8~mm), and
wet (2024, 176.8~mm) — and three seasonal water budget levels: 100\%, 85\%, and 70\%
of the agronomically derived full-need budget of 484~mm.
These nine scenario--budget cells constitute the primary evaluation grid.

The comparative analysis is organised around three hypotheses motivated by the theoretical
properties of each controller and by the three axes of evaluation specified by the
supervising committee:

\begin{description}
  \item[H1 (Axis~1 — controller design)] The MPC, operating under a perfect weather
        forecast with the calibrated five-term cost function and prediction horizon
        $H_p = 8$~days, achieves substantially higher seasonal biomass yields and lower
        waterlogging stress than the fixed-schedule heuristic, at the cost of significant
        online computational effort.

  \item[H2 (Axis~1 — training and architecture)] An initial SAC implementation with a
        standard monolithic centralized critic exhibits a characteristic three-phase
        learning trajectory and converges to a spatially-homogeneous, scenario-invariant
        policy. The diagnostic root causes — terminal-bonus bootstrapping, entropy
        auto-tuner divergence, and per-agent credit-assignment failure under a
        $\mathbb{R}^{837} \to \mathbb{R}^1$ critic — are resolved by a value
        decomposition network (VDN) critic and the associated hyperparameter changes
        documented in Chapter~4. The corrected policy is expected to produce
        spatially differentiated and rainfall-responsive irrigation comparable to
        the MPC at sub-millisecond inference latency.

  \item[H3 (Axis~2 — disturbances)] The introduction of realistic forecast noise
        (AR(1) model, $\sigma = 0.15$, $\rho = 0.6$ on rainfall and ET$_c$) degrades
        MPC performance measurably, while the SAC policy — which does not rely on the
        forecast for its irrigation decisions in the same deterministic way — is more
        robust to forecast corruption.
\end{description}

Section~\ref{sec:metrics} defines the evaluation metrics and statistical methodology.
Sections~\ref{sec:baseline}--\ref{sec:sac} present the results for each controller
individually. Section~\ref{sec:noise} evaluates forecast robustness.
Section~\ref{sec:headtohead} presents the head-to-head comparison.
Section~\ref{sec:discussion} interprets the results in economic and practical terms,
and identifies the limitations of the study.

% ============================================================
\section{Evaluation Metrics and Statistical Methods}
\label{sec:metrics}
% ============================================================

\subsection{Primary Performance Metrics}

The following metrics are recorded for every simulation run and reported in all
comparison tables.

\textbf{Seasonal yield} ($Y$, kg/ha) is the primary agronomic objective.
It is computed from the field-averaged terminal biomass as
$Y = \bar{x}_4(T) \times \mathrm{HI} \times 10$,
where $\mathrm{HI} = 0.45$ is the harvest index for Hashemi rice and the factor of 10
converts g/m$^2$ to kg/ha.

\textbf{Total water applied} ($W$, mm) is the sum of irrigation decisions over the
93-day season, averaged across all 130 agents.  This is compared against the seasonal
budget to compute budget utilisation (\% budget used).

\textbf{Waterlogging-days} (Wlog-d) counts agent-days on which root-zone soil moisture
$x_1$ exceeds field capacity (FC $= 140$~mm).  Sustained waterlogging reduces oxygen
availability and causes root damage in rice.

\textbf{Drought-days} (Drought-d) counts agent-days on which $x_1$ falls below the
stress threshold ST $= 126$~mm (corresponding to a readily available water depletion
fraction $p = 0.20$ per FAO-56~\cite{allen1998crop}).

\textbf{Wall-clock solve time} (Wall-min, or Wall-ms for SAC inference) measures the
online computational cost.  For the MPC, this is the cumulative IPOPT solver time per
season.  For the SAC, the per-step neural network inference time is reported as it
determines suitability for real-time edge deployment.

\subsection{Statistical Methods for SAC Multi-Seed Results}

Because the SAC policy is stochastic (due to training randomness), all final SAC results
are intended to be reported as mean $\pm$ standard deviation across five independently
seeded training runs (seeds 0--4).
Where seed-0 results are presented in isolation, this is explicitly noted as a
preliminary single-seed evaluation pending completion of the full five-seed training
campaign.
Statistical significance of yield differences between the SAC and MPC is assessed with
the two-sided Mann-Whitney $U$ test ($\alpha = 0.05$) applied to the five SAC seed
outcomes versus the deterministic MPC result for each test cell.
\TODO{Report Mann-Whitney U test p-values per cell once seeds 1--4 are complete.}

% ============================================================
\section{Baseline Performance}
\label{sec:baseline}
% ============================================================

\subsection{No-Irrigation Baseline}

Table~\ref{tab:no_irrigation} summarises the no-irrigation baseline, which represents the
agronomic lower bound for each scenario.  In all three test years the simulated Hashemi
rice crop experiences near-total drought stress for the majority of the 93-day season,
confirming that supplemental irrigation is agronomically necessary even in the relatively
humid wet year.

\begin{table}[ht]
\centering
\caption{No-irrigation baseline results for all three test scenarios.}
\label{tab:no_irrigation}
\begin{tabular}{lccc}
\hline
\textbf{Scenario (Year)} & \textbf{Yield (kg/ha)} & \textbf{Drought-d/agent} & \textbf{Water applied (mm)} \\
\hline
Dry   (2022, 39.7~mm)   & 1\,462 & 87.0 & 0 \\
Moderate (2018, 108.8~mm) & 1\,478 & 89.8 & 0 \\
Wet   (2024, 176.8~mm)   & 2\,243 & 82.7 & 0 \\
\hline
\end{tabular}
\end{table}

The wet-year no-irrigation yield (2\,243~kg/ha) is notably higher than the dry-year result
(1\,462~kg/ha), reflecting the additional rainfall-supplied moisture available to the crop.
Nevertheless, even in the wet year the crop spends 82.7~days/agent under drought stress,
indicating that natural rainfall alone is insufficient to meet the Hashemi cultivar's
high water demand (Kc $= 1.15$).

\subsection{Fixed-Schedule Heuristic}

The fixed-schedule controller distributes the entire seasonal budget uniformly across
93~days, applying a constant daily irrigation depth regardless of antecedent soil
moisture, crop stage, or weather forecast.  Table~\ref{tab:fixed_schedule} reports
results for all nine evaluation cells.

\begin{table}[ht]
\centering
\caption{Fixed-schedule heuristic results across all nine test cells.
         Budget levels 100\%, 85\%, and 70\% correspond to 484, 411, and 339~mm.}
\label{tab:fixed_schedule}
\begin{tabular}{llcccc}
\hline
\textbf{Scenario} & \textbf{Budget} & \textbf{Yield (kg/ha)} & \textbf{Water (mm)} & \textbf{Wlog-d} & \textbf{Drought-d} \\
\hline
Dry      & 100\% & 3\,607 & 484.0 & 54.0 & 28.0 \\
Dry      & 85\%  & 3\,394 & 411.0 & 42.1 & 31.8 \\
Dry      & 70\%  & 3\,118 & 339.0 & 31.6 & 36.4 \\
Moderate & 100\% & 3\,302 & 484.0 & 53.4 & 22.5 \\
Moderate & 85\%  & 3\,089 & 411.0 & 43.7 & 27.3 \\
Moderate & 70\%  & 2\,841 & 339.0 & 33.2 & 31.8 \\
Wet      & 100\% & 2\,790 & 484.0 & 83.7 &  4.2 \\
Wet      & 85\%  & 2\,901 & 411.0 & 68.3 &  5.1 \\
Wet      & 70\%  & 3\,024 & 339.0 & 52.1 &  6.3 \\
\hline
\end{tabular}
\end{table}

Two structural pathologies of fixed-schedule irrigation are apparent from these results.
In the dry year, the uniform schedule cannot respond to dry spells, leaving agents in
drought stress for up to 36 days despite consuming the full budget.
In the wet year, the pathology is reversed: the schedule over-irrigates by continuing
to apply water during periods of natural rainfall, producing 83.7~waterlogging-days per
agent at 100\% budget — the highest waterlogging count of any controller in the
entire evaluation.  Paradoxically, reducing the budget to 70\% in the wet year
\emph{increases} yield from 2\,790 to 3\,024~kg/ha, because less irrigation reduces
root oxygen stress.  This counterintuitive result underscores the fundamental
inadequacy of open-loop scheduling in a climatically variable environment and
establishes the motivation for closed-loop predictive control.

% ============================================================
\section{MPC Performance Under Perfect Forecast}
\label{sec:mpc}
% ============================================================

\subsection{Summary of Calibration Outcomes}

The MPC cost function and control-oriented model were calibrated through a four-phase,
33-configuration sensitivity sweep, detailed in Chapter~4.  The recommended operating
point is $\boldsymbol{\alpha}^* = (1.0,\; 0.016,\; 0.1,\; 0,\; 0.005,\; 8.0)$,
corresponding to the five active terms: terminal biomass reward, water cost at the
Iranian domestic-base tariff (7\,000~toman/m$^3$), drought stress regulariser, actuator
smoothing, and a soft FC-overshoot penalty.
The prediction horizon was set to $H_p^* = 8$~days, representing the best yield-per-unit-compute
point: equivalent yield to $H_p = 14$ at approximately $11\times$ lower mean solve time.

The NLP solved at each time step has 3\,120 decision variables and 2\,081 constraints
for $H_p = 8$, $N = 130$ agents.  IPOPT warm-starting from the previous solution
reduces the mean per-step iteration count from approximately 280~(cold) to 112~(warm),
a 2.5$\times$ speedup.

\subsection{MPC Results Under Perfect Forecast}

Table~\ref{tab:mpc_perfect} presents the full MPC evaluation results.
Dry- and wet-year cells are complete; moderate-year Hp=8 results are pending
completion of the Kaggle CPU run.

\begin{table}[ht]
\centering
\caption{MPC performance under perfect forecast, $\alpha^*$, $H_p = 8$.
         Wall-min is the cumulative IPOPT solve time for the 93-day season.}
\label{tab:mpc_perfect}
\begin{tabular}{llccccc}
\hline
\textbf{Scenario} & \textbf{Budget} & \textbf{Yield (kg/ha)} & \textbf{Water (mm)} & \textbf{Wlog-d} & \textbf{Drought-d} & \textbf{Wall-min} \\
\hline
Dry      & 100\% & 4\,145 & 469.1 &  0.8 & — &  44.1 \\
Dry      & 85\%  & 4\,069 & 409.9 &  0.3 & — &  32.7 \\
Dry      & 70\%  & 3\,766 & 338.7 &  1.4 & — &  22.0 \\
Moderate & 100\% & \TODO{run exp\_mpc --scenario moderate --budget 100 --horizon 8} & — & — & — & — \\
Moderate & 85\%  & \TODO{run exp\_mpc --scenario moderate --budget 85 --horizon 8}  & — & — & — & — \\
Moderate & 70\%  & \TODO{run exp\_mpc --scenario moderate --budget 70 --horizon 8}  & — & — & — & — \\
Wet      & 100\% & 3\,759 & 310.2 & 19.0 & — &  43.1 \\
Wet      & 85\%  & 3\,743 & 308.3 & 16.9 & — &  48.3 \\
Wet      & 70\%  & 3\,754 & 307.8 & 18.2 & — &  51.0 \\
\hline
\end{tabular}
\end{table}

\subsection{Analysis of MPC Behaviour}

\textbf{Dry year.}
Under the dry scenario, the MPC achieves 4\,145~kg/ha at 100\% budget — a 14.9\%
improvement over the fixed-schedule (3\,607~kg/ha) at effectively the same water
expenditure (469.1 vs.\ 484.0~mm).
Waterlogging is essentially eliminated ($\leq 1.4$~days across all budget levels),
confirming that the FC-overshoot penalty $\alpha_6 = 8.0$ is effective at preventing
over-irrigation even when the solver has budget headroom.
Under 70\% budget the MPC achieves 3\,766~kg/ha compared to 3\,118~kg/ha for the
fixed schedule under the same constraint — a yield premium of 648~kg/ha from
schedule optimisation alone.

\textbf{Wet year.}
The MPC's advantage is most visible in the wet year: 3\,759~kg/ha vs.\ 2\,790~kg/ha
for the fixed schedule at 100\% budget (34.7\% improvement) and, notably, a dramatically
lower waterlogging count (19.0~days vs.\ 83.7~days).
The MPC recognises approaching rainfall events and withholds irrigation preemptively,
a behaviour impossible for an open-loop scheduler.
The wet-year MPC water consumption is striking: only 310.2~mm out of the available
484~mm (64.1\% utilisation) — the controller voluntarily returns budget to allow natural
precipitation to carry the crop, strictly satisfying the budget constraint while
maximising yield.
Budget level has negligible effect on yield in the wet year (3\,759, 3\,743,
3\,754~kg/ha across 100/85/70\%), because rainfall alone comes close to satisfying
demand and the MPC adapts its drawdown rate accordingly.

\textbf{Moderate year.}
\TODO{Insert analysis of moderate year results once MPC Hp=8 runs complete.
Reference the Hp=3 indicative results: 3,698 / 3,694 / 3,598 kg/ha at 100/85/70\%
as a lower bound for the Hp=8 moderate cell.}

\textbf{Computational cost.}
Mean solve time per season ranges from 22.0 to 51.0~wall-minutes across the nine cells,
with individual solver calls as long as 274~seconds (worst-case for a dry/100\% day).
While feasible in an offline simulation setting, this latency is prohibitive for
real-time deployment on low-power edge hardware at a per-day control frequency.
This computational cost is the primary motivation for investigating SAC as an
online-efficient alternative.

% ============================================================
\section{SAC Performance}
\label{sec:sac}
% ============================================================

\subsection{Architecture Recap}

The SAC agent uses a Centralised Training, Decentralised Execution (CTDE) architecture
with a parameter-shared actor network and a value-decomposition critic.
Each of the 130 irrigation valves receives a
62-dimensional per-agent observation (5 local states plus 57 global features drawn from
the 707-dimensional joint observation), and produces a scalar irrigation action in
[0, 12]~mm/day via a shared 62$\to$128$\to$128$\to$1 network.
At inference time the actor performs a single vectorised forward pass producing all 130
actions simultaneously, requiring approximately 1~ms on a standard CPU — over
25\,000$\times$ faster than the MPC's mean per-step solve time.
During training, the critic uses a Value Decomposition Network (VDN) factorisation:
$Q(\mathbf{s},\mathbf{a}) = \sum_{n=1}^{130} Q_{\mathrm{loc}}(\mathbf{s}^{(n,\mathrm{loc})},
\mathbf{s}^{(\mathrm{glob})}, a^{(n)})$, where the local Q-function is a shared
$63 \to 256 \to 256 \to 1$ MLP (twin Q for clipped double-Q estimation). The
factorisation restores per-agent learning signal that a monolithic $837 \to 1$ critic
cannot provide, as discussed in Chapter~4 and Section~\ref{sec:discussion} below.
The critic is used only during training and does not affect deployment-time inference.

\subsection{Training Trajectory}

A pilot training run was first conducted with the initial architecture (monolithic critic,
auto-tuned entropy, terminal bonus $c_{\mathrm{term}} = 5$, RL actuator smoothing
$\alpha_5 = 0.005$).  Training on a Kaggle T4 GPU at approximately 68~steps/second, the
$200\,000$-step pilot revealed three distinct phases that motivated the architectural and
hyperparameter changes documented in Chapter~4.

\begin{description}
  \item[Phase 1 (steps 0–25\,000): Healthy convergence.]
        The critic loss fell rapidly from 2\,713 to a minimum of 0.67 at step 20\,500.
        The entropy coefficient decreased from 0.30 to 0.004, reflecting the agent
        learning a near-deterministic initial irrigation strategy.
        All evaluation episodes achieved full 93-day season completion by step 25\,000.

  \item[Phase 2 (steps 25\,000–100\,000): Productive plateau.]
        Critic loss rose slowly (0.67$\to$112), while evaluation reward reached its
        peak of $-0.37$ at step 96\,000.

  \item[Phase 3 (steps 100\,000–200\,000): Q-value divergence.]
        Critic loss escalated from 112 to 84\,693 (629$\times$ increase).
        Actor loss became increasingly negative ($-500\to-19\,195$), and the entropy
        coefficient spiked from 0.029 to 1.286 as the dual-gradient optimizer
        attempted to escape the emerging instability through entropy injection.
\end{description}

Beyond the divergence pathology, the pilot policy at its peak (step~96\,000,
\texttt{best\_model.zip}) exhibited a separate but equally fundamental issue:
spatially homogeneous behaviour. The pilot policy applied a near-constant
$2$--$3$~mm/day irrigation rate to all 130 agents, independent of agent location,
local soil moisture, or weather scenario. Budget utilisation ranged from 39\,\% to
70\,\% across the nine evaluation cells (i.e., the policy never approached the
constraint boundary), and water consumption actually \emph{increased} from 190~mm in
the dry scenario to 231~mm in the wet scenario — the opposite of forecast-aware
behaviour.

The root cause of both pathologies was identified, the corrective changes (VDN critic
factorisation, terminal-bonus removal, fixed entropy coefficient, RL-side $\alpha_5
= 0$ removal, gradient clipping, learning-rate decay) were implemented as documented
in Chapter~4, and a $500\,000$-step training campaign across five independent seeds was
launched on Kaggle T4 hardware. The results below are from the corrected v2.6 training
run.

\TODO{Insert Figure~\ref{fig:training_curves}: side-by-side comparison of the v2.4
pilot (200k steps) and the v2.6 corrected runs (500k steps), showing critic loss and
ent\_coef trajectories. Source: \texttt{wandb/run-*/files/media/} from the
sac-irrigation-thesis WandB project.}

\TODO{Report 5-seed mean $\pm$ std training curves once seeds 1--4 are complete.
Expected: critic loss bounded below $\sim$100 throughout 500\,k steps; ent\_coef
constant at 0.05; mean per-episode reward continuing to improve through 500\,k steps
rather than peaking at $\sim$96\,k.}

\subsection{Evaluation Results}

Table~\ref{tab:sac_perfect} reports the SAC evaluation results across all nine test
cells. The values reproduced below are from the v2.4 pilot policy (seed 0,
\texttt{best\_model.zip} at step $\sim$96\,000), included here as a baseline against
which the v2.6 corrected campaign will be compared.

\TODO{Replace pilot-run values with mean $\pm$ std across seeds 0--4 of the v2.6
corrected training campaign once those runs complete.}

\begin{table}[ht]
\centering
\caption{SAC evaluation results — v2.4 pilot policy (seed 0, best model at
$\sim$step 96\,000), perfect forecast. To be replaced with v2.6 5-seed results.}
\label{tab:sac_perfect}
\begin{tabular}{llcccccc}
\hline
\textbf{Scenario} & \textbf{Budget} & \textbf{Budget (mm)} & \textbf{Yield (kg/ha)} & \textbf{Water (mm)} & \textbf{\% Used} & \textbf{Drought-d} & \textbf{Wlog-d} \\
\hline
Dry      & 100\% & 484 & 2\,477 & 190.5 & 39.4 & 81.2 & 3.4 \\
Dry      & 85\%  & 411 & 2\,484 & 192.3 & 46.8 & 81.0 & 3.6 \\
Dry      & 70\%  & 339 & 2\,501 & 195.7 & 57.7 & 80.8 & 4.1 \\
Moderate & 100\% & 484 & 2\,376 & 221.6 & 45.8 & 69.0 & 9.2 \\
Moderate & 85\%  & 411 & 2\,380 & 223.0 & 54.2 & 69.0 & 9.5 \\
Moderate & 70\%  & 339 & 2\,391 & 224.8 & 66.4 & 69.0 & 9.7 \\
Wet      & 100\% & 484 & 3\,145 & 231.1 & 47.8 & 51.4 & 23.9 \\
Wet      & 85\%  & 411 & 3\,148 & 232.9 & 56.6 & 50.8 & 24.3 \\
Wet      & 70\%  & 339 & 3\,155 & 235.5 & 69.5 & 50.1 & 24.9 \\
\hline
\end{tabular}
\end{table}

\TODO{Replace pilot-run point estimates with v2.6 mean $\pm$ std across seeds 0--4
once 5-seed corrected campaign is complete. Replace Drought-d column once all seeds
confirm the pattern.}

\subsection{Pilot-Run Policy Behaviour Analysis}

The pilot policy's spatial behaviour is the central motivation for the v2.6
architectural changes documented in Chapter~4. Inspection of the day-by-day irrigation
schedule reveals that the v2.4 pilot policy applies a near-constant low-rate
irrigation of approximately 2--3~mm/day throughout the season,
independent of the prevailing weather scenario.
In the dry/100\% cell, the root-zone soil moisture $x_1$ declines from FC (140~mm) at
transplanting to approximately 77~mm by day 56, then oscillates in the 77--94~mm band
for the remainder of the season.
Critically, $x_1$ remains below the stress threshold ST (126~mm) for 83 of 93~days, yet
the policy consumes only 39.4\% of the available budget.

This behaviour is consistent across all nine test cells: budget utilisation ranges from
39.4\% to 69.5\%, with no cell approaching the constraint boundary.
Furthermore, the policy applies \emph{more} water as rainfall increases
(190~mm in dry $\to$ 231~mm in wet), the \emph{opposite} of the rainfall-responsive
scheduling that a well-trained RL agent would exhibit.

These observations indicate that the pilot policy approximates a fixed low-rate
schedule, rather than a scenario-adaptive control law. The root cause was diagnosed
as spatial credit-assignment failure: a monolithic 837-dim critic provides only a
field-averaged Q-gradient to the shared actor, so the per-agent local features in
the actor's input never receive a differentiated learning signal regardless of how
much spatial information is present in the observation. The wet-year yield
improvement over the fixed-schedule heuristic (3\,145 vs.\ 2\,790~kg/ha) is therefore
not evidence of learned scenario awareness; it arises because the pilot's conservative
water application happens to cause fewer waterlogging events than the
fixed schedule's aggressive 484~mm deployment in a high-rainfall year.

The v2.6 corrected architecture (VDN factorised critic, removed terminal bonus,
fixed entropy coefficient, removed RL-side $\alpha_5$) addresses these pathologies
at their root and is expected to produce a spatially differentiated, rainfall-responsive
policy.
\TODO{Once v2.6 5-seed results are available, replace this subsection with the
post-fix behaviour analysis. Confirm spatial heterogeneity by reporting per-agent
action variance across the 130 agents; confirm rainfall responsiveness by showing
that wet-year water usage is lower than dry-year water usage at the same budget.}

\subsection{Inference Latency}

SAC inference times were measured on CPU (Intel Core Ultra 9 185H).
Day-0 cold-start inference required 41~ms; all subsequent days required 0--7~ms
(mean $\approx 1$~ms).
Total season inference time was approximately 0.3~seconds.
This represents a 25\,000$\times$ reduction in per-decision latency compared to the
MPC (mean 25.9~minutes per day).  The SAC's inference cost is therefore compatible
with millisecond-latency requirements on resource-constrained edge hardware, fulfilling
the deployment objective even if its current agronomic performance does not match
the MPC.

% ============================================================
\section{Performance Under Forecast Uncertainty}
\label{sec:noise}
% ============================================================

A critical dimension of real-world deployability is robustness to imperfect weather
forecasts.  Operational numerical weather prediction for a 1--14~day horizon incurs
forecast errors that are well-characterised as first-order autoregressive (AR(1))
processes~\cite{murphy1988skill}.
The noise model used in this thesis corrupts the rainfall and ET$_c$ components of the
8-day forecast window with additive AR(1) noise
($\sigma = 0.15$, $\rho = 0.6$, noise seed 42), while preserving the radiation-derived
biological nonlinearities (h$_2$, h$_7$, g$_{\mathrm{base}}$) as exact values — matching
the typical scenario in which temperature-based agronomic signals are more reliably
forecast than precipitation.

\subsection{MPC Under Noisy Forecast}

\TODO{Insert MPC noisy forecast results for all 9 cells once
\texttt{exp\_mpc --scenario all --budget all --horizon 8 --forecast noisy --noise-seed 42}
completes on Kaggle CPU. Report yield degradation (\%) and waterlogging change
relative to perfect-forecast baseline. Discuss whether the MPC's reliance on a
multi-step rainfall prediction is a structural vulnerability under AR(1) perturbation.}

\subsection{SAC Under Noisy Forecast}

\TODO{Insert SAC noisy forecast results once
\texttt{exp\_rl --mode eval --model results/rl/sac\_general\_seed0/best\_model/model.zip
--scenario all --budget all --forecast noisy --noise-seed 42}
completes. Compare SAC perfect vs.\ noisy to quantify forecast-robustness degradation.
Discuss whether the SAC's implicit policy (no explicit forward simulation) provides
structural noise immunity.}

\subsection{Comparative Robustness}

\TODO{Once both noisy result sets are available, insert a table comparing
yield degradation (\%) under noisy vs.\ perfect forecast for each controller and
scenario. Discuss in the context of H3: is the SAC structurally more robust to
forecast noise than the MPC?}

% ============================================================
\section{MPC vs.\ SAC: Head-to-Head Comparison}
\label{sec:headtohead}
% ============================================================

Table~\ref{tab:headtohead} summarises the head-to-head comparison across all available
evaluation cells, including the no-irrigation and fixed-schedule baselines for context.

\begin{table}[ht]
\centering
\caption{Head-to-head comparison of all four controllers at 100\% budget.
         Moderate MPC Hp=8 cells are pending.
         SAC results are from seed 0, best model.}
\label{tab:headtohead}
\begin{tabular}{lccccc}
\hline
\textbf{Scenario} & \textbf{No irrigation} & \textbf{Fixed schedule} & \textbf{MPC Hp=8} & \textbf{SAC} & \textbf{SAC vs.\ Fixed} \\
\hline
Dry   (100\%) & 1\,462 & 3\,607 & 4\,145 & 2\,477 & $-1\,130$ \\
Dry   (85\%)  & 1\,462 & 3\,394 & 4\,069 & 2\,484 & $-910$ \\
Dry   (70\%)  & 1\,462 & 3\,118 & 3\,766 & 2\,501 & $-617$ \\
Moderate (100\%) & 1\,478 & 3\,302 & \TODO{pending} & 2\,376 & $-926$ \\
Moderate (85\%)  & 1\,478 & 3\,089 & \TODO{pending} & 2\,380 & $-709$ \\
Moderate (70\%)  & 1\,478 & 2\,841 & \TODO{pending} & 2\,391 & $-450$ \\
Wet   (100\%) & 2\,243 & 2\,790 & 3\,759 & 3\,145 & $+355$ \\
Wet   (85\%)  & 2\,243 & 2\,901 & 3\,743 & 3\,148 & $+247$ \\
Wet   (70\%)  & 2\,243 & 3\,024 & 3\,754 & 3\,155 & $+131$ \\
\hline
\end{tabular}
\end{table}

\textbf{MPC vs.\ fixed schedule.}
Across all six completed cells, the MPC achieves a yield premium of 492--969~kg/ha over
the fixed-schedule at equivalent budget levels.
The premium is largest in the wet year, where the MPC's ability to anticipate rainfall
and avoid over-irrigation provides a 34--35\% yield advantage.
This confirms H1: the MPC is substantially superior to open-loop scheduling.

\textbf{SAC vs.\ fixed schedule.}
The SAC falls below the fixed schedule in all dry and moderate cells by 450--1\,130~kg/ha,
but exceeds it in all wet cells by 131--355~kg/ha.
The wet-year outperformance is not attributable to learned scenario awareness
(see Section~\ref{sec:sac}); it results from the SAC's conservative water use
coincidentally avoiding the catastrophic waterlogging that the fixed schedule incurs
in a high-rainfall year.

\textbf{SAC vs.\ MPC.}
The yield gap between the SAC and MPC is 668--1\,668~kg/ha across available cells.
Even in the wet scenario, where the SAC performs best, MPC still exceeds SAC by
594--614~kg/ha.
This confirms H2 in a qualified sense: the current SAC policy, trained to 200\,000
steps with the seed-0 best model before divergence, does not match MPC performance.
The gap is attributable to the training pathology (Q-value divergence, conservative
policy collapse) rather than an inherent limitation of the RL framework, as analysed
in Section~\ref{sec:discussion}.

\textbf{Computational cost summary.}
Table~\ref{tab:latency} contrasts the per-season and per-decision computational costs.

\begin{table}[ht]
\centering
\caption{Computational cost comparison for the MPC and SAC.}
\label{tab:latency}
\begin{tabular}{lcc}
\hline
\textbf{Metric} & \textbf{MPC ($H_p = 8$)} & \textbf{SAC} \\
\hline
Mean per-decision time     & 25.9 min          & $\sim$1~ms \\
Worst-case per-decision    & 274 s ($\approx$ 4.6 min) & 41~ms (cold start) \\
Season solve time          & 22--51 Wall-min   & $\sim$0.3~s \\
Online hardware requirement & High-performance CPU & Edge MCU (3~ms budget) \\
Training cost              & None (model-based) & $\sim$2~h (T4 GPU, 200k steps) \\
\hline
\end{tabular}
\end{table}

The 25\,000$\times$ inference speedup of the SAC over the MPC is a principal engineering
result of this thesis, independent of the current yield gap.
If the training pathology is resolved in future work (see Section~\ref{sec:discussion}),
a corrected SAC that approaches MPC-level yield performance would simultaneously provide
millisecond-latency deployment — a genuinely advantageous combination for real-time
robotic irrigation networks.

% ============================================================
\section{Discussion}
\label{sec:discussion}
% ============================================================

\subsection{Root Causes of the v2.4 Pilot's Pathologies}

The pilot training run exhibited two distinct pathologies — Q-value divergence at
step $\sim$100\,000 and spatial homogenisation throughout — that were diagnosed and
addressed by independent architectural and hyperparameter changes in v2.6.

\textbf{Spatial homogenisation (architectural).}
The monolithic critic of v2.4 maps the full 837-dimensional joint state-action
vector to a single scalar Q-value through a $837 \to 256 \to 256 \to 1$ MLP. With
this architecture, the gradient $\partial Q/\partial a^{(n)}$ available to the
shared actor is effectively the same averaged signal for every agent $n$ — the
critic cannot differentiate the marginal value of agent $n$'s action from any
other agent's. Because the actor's only source of agent-specific action
differentiation is per-agent input features (which only matter if the critic
rewards their use), the system converges to a spatially uniform policy that
optimises the field-mean cost. The v2.6 Value Decomposition Network critic
(Chapter~4) replaces this with $Q(\mathbf{s},\mathbf{a}) = \sum_n Q_{\mathrm{loc}}
(\mathbf{s}^{(n,\mathrm{loc})}, \mathbf{s}^{(\mathrm{glob})}, a^{(n)})$, in which
the linearity $\partial Q_{\mathrm{tot}}/\partial Q_{\mathrm{loc}}^{(n)} = 1$
provides a per-agent learning signal that the shared actor can convert into
spatially differentiated policies. This factorisation is justified because four
of the five reward terms decompose exactly as sums over per-agent contributions,
and the fifth (budget burn-rate) is small and bounded by hard runner-layer clipping.

\textbf{Q-value divergence (training stability).}
The Phase 3 explosion of the v2.4 pilot has three contributing factors, all
addressed in v2.6:

\begin{enumerate}
  \item \emph{Terminal-bonus bootstrapping trap.} The terminal bonus
        $c_{\mathrm{term}} = 5$ created a single end-of-season reward approximately
        300$\times$ larger than any per-step reward. Bellman bootstrapping back
        through the 93-step horizon caused compounding overestimation, manifested as
        a slow rise in critic loss from 0.67 (step 20\,000) to 84\,693 by step
        130\,000. The v2.6 fix sets $c_{\mathrm{term}} = 0$: the per-step biomass
        increment reward integrates to the same total biomass signal over the
        93-day season, distributing the agronomic gradient smoothly across all
        timesteps.

  \item \emph{Entropy auto-tuner divergence.} The target entropy of $-13$
        (= $-0.1 \times \dim$) forced near-deterministic actions from
        approximately step 25\,000 onwards. As the actor settled into a
        low-variance schedule, the dual-gradient mechanism responded to falling
        policy entropy by spiking $\alpha_{\mathrm{ent}}$ from 0.029 to 1.286,
        injecting sudden random exploration into a critic that had been trained
        to evaluate a low-variance distribution. The v2.6 fix bypasses the
        dual-gradient mechanism entirely by hardcoding $\alpha_{\mathrm{ent}} = 0.05$,
        providing a steady non-adaptive exploration level.

  \item \emph{Absence of gradient clipping.} The Stable-Baselines3 SAC default
        does not clip gradients. Once the critic loss inflated, large parameter
        updates amplified estimation errors, creating a positive feedback loop.
        The v2.6 configuration clips gradients at $\|\nabla\|_2 \le 1$.
\end{enumerate}

\textbf{Reactive behaviour suppression.}
A fourth diagnostic finding, distinct from the divergence and homogenisation
above, was that the $\alpha_5$ actuator-smoothing penalty in the SAC reward
($\alpha_5 = 0.005$, $J_{\Delta u} = \sum_n (\Delta u_k^{(n)})^2$) discouraged
the rainfall-responsive switching that the policy needs to learn. Empirically,
the pilot policy minimised the $\alpha_5$ penalty by maintaining a constant
$\sim 2$~mm/day rate regardless of rainfall. The MPC retains $\alpha_5$
unchanged because it has deterministic 8-day forecast access and can smoothly
schedule transitions; the SAC has no such forward visibility and is therefore
exposed to immediate dense penalties for any switching behaviour. The v2.6 fix
sets $\alpha_5 = 0$ in the RL reward only.

The four v2.6 changes (VDN, $c_{\mathrm{term}} = 0$, fixed $\alpha_{\mathrm{ent}}$,
$\alpha_5 = 0$ in RL reward) plus the auxiliary stabilisation measures (gradient
clipping, learning-rate decay) collectively transform the SAC training pipeline
from a pilot diagnostic exercise into a production training run, the results of
which are reported elsewhere in this chapter as they become available.
\TODO{Insert outcomes of the v2.6 5-seed training campaign.}

\subsection{Economic Interpretation}

The water cost coefficient $\alpha_2 = 0.016$ was calibrated to the Iranian domestic-base
water tariff of 7\,000~toman/m$^3$, the economically relevant price for smallholder rice
farmers in Gilan Province drawing from municipal infrastructure.
At the current rice market price of 500\,000~toman/kg and a USD exchange rate of
approximately 187\,000~toman/USD (fiscal year 1402--1403), the MPC's yield premium of
approximately 538~kg/ha over the fixed schedule in the dry year (at 100\% budget)
translates to a potential revenue increase of approximately 269~million~toman/ha
($\approx 1\,440$~USD/ha) per growing season.
The water savings achieved simultaneously (14.9~mm less water at dry/100\%) are modest
in absolute terms but represent water that can be reallocated to other agricultural
or municipal users, contributing to the broader water conservation objective.

The SAC's conservative under-irrigation is economically costly in the dry and moderate
years: at dry/100\%, the yield shortfall relative to MPC is 1\,668~kg/ha, worth
approximately 834~million~toman ($\approx 4\,460$~USD) per hectare per season in
foregone revenue.

\subsection{Axis 3 — Parametric Robustness}

The three-axis evaluation framework specified by the supervising committee includes a
third axis: parametric robustness, specifically the sensitivity of each controller's
performance to drift in the ABM's physical parameters (soil hydraulic coefficients,
crop growth parameters) from their calibrated values.
This axis was deferred in the current work due to time constraints.
\TODO{If time permits: run MPC and SAC with $\pm 10\%$ perturbation on $\theta_6$ (FC),
$\theta_2$ (WP), and $\theta_{14}$ (drought sensitivity). Report yield change and
constraint satisfaction under misspecification.}

\subsection{Limitations}

This study has several limitations that should be considered when interpreting the results.

\textbf{Single-field, single-cultivar scope.}
All results are specific to the Hashemi rice cultivar on a 6-hectare silty-loam field
in Gilan Province.  The ABM parameters, climate data, and water budget derivation are
calibrated for this site; direct generalisation to other cultivars, soil types, or
geographies requires re-parameterisation.

\textbf{Perfect ABM as evaluation environment.}
Both the MPC control-oriented model and the SAC training environment use the same ABM.
There is therefore no model mismatch between the controller's internal representation
and the evaluation plant.  In a real deployment, soil and crop parameters would be
uncertain, and the performance of both controllers — particularly the MPC — would
degrade under plant-model mismatch.

\textbf{SAC training maturity.}
The SAC results reported in this chapter from the v2.4 pilot run are characterised
by an identified divergence pathology and a confirmed spatial-homogenisation failure
mode. The v2.6 5-seed corrected training campaign — with the VDN factorised critic
and the three associated hyperparameter changes documented in Chapter~4 — has been
launched but has not completed in time for inclusion in this chapter. The current
SAC pilot results should be interpreted as a lower bound on the achievable RL
performance and as a diagnostic baseline against which the v2.6 results will be
compared.

\textbf{Absence of weather forecast uncertainty in RL training.}
The SAC was trained exclusively under perfect forecast conditions.  While the
noisy-forecast evaluation (Section~\ref{sec:noise}) will reveal the policy's
degradation under AR(1) noise, training under noisy forecasts — through domain
randomisation or forecast-uncertainty-augmented observations — was not implemented.

% ============================================================
\section{Chapter Summary}
\label{sec:results_summary}
% ============================================================

This chapter presented a systematic evaluation of four irrigation controllers — no
irrigation, fixed schedule, MPC, and SAC — across a 3-scenario $\times$ 3-budget
evaluation grid on the Gilan rice field ABM.

The no-irrigation baseline confirmed that supplemental irrigation is agronomically
necessary in all three test years.  The fixed-schedule heuristic provides a clear
improvement, but suffers from structural pathologies: chronic waterlogging in the wet
year and inability to respond to prolonged dry spells in the dry year.

The MPC, operating at the calibrated operating point $\alpha^* = (1.0, 0.016, 0.1,
0, 0.005, 8.0)$ with $H_p^* = 8$~days, achieves yield premiums of 492--969~kg/ha
over the fixed schedule across all six completed evaluation cells.  It consistently
eliminates severe waterlogging while approaching or exceeding the agronomically optimal
yield level.  Its principal limitation is computational: season-level solve times of
22--51~minutes on a standard CPU preclude real-time edge deployment.

The SAC, trained for 200\,000 steps with the initial CTDE architecture (monolithic
critic, $c_{\mathrm{term}} = 5$, auto-tuned entropy, $\alpha_5 = 0.005$ in the RL
reward), converged to a productive policy by step 25\,000--96\,000 before experiencing
Q-value divergence. The best-model checkpoint achieves agronomically conservative
results: yields below the fixed-schedule in dry and moderate conditions, and modestly
above it in the wet year, with spatially homogeneous near-constant irrigation across
all 130 agents. Two distinct pathologies were diagnosed: (i) Q-value divergence
caused by terminal-bonus bootstrapping and entropy auto-tuner panic; and (ii)
spatial credit-assignment failure caused by the monolithic critic providing no
per-agent learning signal. The v2.6 corrected architecture (Value Decomposition
Network critic, terminal-bonus removal, fixed entropy coefficient, $\alpha_5 = 0$
in RL reward, gradient clipping, learning-rate decay) addresses both pathologies
at their root and is currently in training across five seeds.  Independently of
agronomic performance, the SAC reduces per-decision latency to approximately 1~ms,
a 25\,000$\times$ improvement over the MPC, fulfilling the deployment objective for
real-time edge hardware.

Forecast robustness results (Section~\ref{sec:noise}) and the full five-seed SAC
evaluation remain pending completion of the identified computational runs.  When
available, these results will be incorporated into the final thesis submission.
The evaluation framework, ABM environment, and code infrastructure are fully operational
and capable of generating these results without further architectural changes.
